\documentclass[10pt,a4paper]{article}
\usepackage{fontspec,amsmath,amssymb,graphicx}
\usepackage[margin=18mm]{geometry}
\setmainfont{Times New Roman}
\setlength{\parindent}{0pt}
\setlength{\parskip}{5pt}
\allowdisplaybreaks
\sloppy
\emergencystretch=2em
\begin{document}
\section*{VD\_3\_2\_2\_01 — Lời giải bổ sung}
\subsection*{Kiến thức và trình tự giải}
\textbf{Bước 1 — Liệt kê các khoảng của dữ liệu.} Ghi đủ điểm đầu, điểm cuối hỗ trợ, điểm gãy và bước nhảy của $f,g$. Nếu có biên, phản xạ chẵn hoặc lẻ trước khi tích phân; trên đoạn hữu hạn tiếp tục phản xạ đến một chu kỳ hoàn chỉnh.

\textbf{Bước 2 — Tính nguyên hàm liên tục.} Lấy $H'=\widetilde g$ trên từng khoảng và chọn hằng số để hai vết của $H$ bằng nhau. Viết $[s]_p^q=\min(q,\max(p,s))$ hoặc $(s)_+=\max(s,0)$ khi cần gộp nhiều nhánh. Một nguyên hàm liên tục tránh tạo ra bước nhảy giả trong nghiệm.

\textbf{Bước 3 — Xác định sóng phải và sóng trái.} Đặt $A(s)=\widetilde f(s)/2-H(s)/(2a)$, $B(s)=\widetilde f(s)/2+H(s)/(2a)$; khi đó $u=A(x-at)+B(x+at)$. Với từng $t$ được yêu cầu, dịch đồ thị $A$ sang phải $at$ và $B$ sang trái $at$, rồi cộng tung độ. Các đường $x-at=b$ và $x+at=b$ chia mặt phẳng thành các miền có công thức đơn giản.

\textbf{Bước 4 — Xét liên tục, đạo hàm và giao của các mặt sóng.} Nếu $f$ có bước nhảy, $u$ thường có bước nhảy; nếu $f$ liên tục và $g$ bị chặn thì $u$ liên tục. Tại thời điểm cố định,


\[
u_x=\frac{\widetilde f'(x-at)+\widetilde f'(x+at)}2+\frac{\widetilde g(x+at)-\widetilde g(x-at)}{2a}.
\]


Tính giới hạn trái/phải của biểu thức này tại từng vị trí ứng viên. Có thể có triệt tiêu khi hai đường gặp nhau, nên không chỉ liệt kê đường rồi khẳng định mọi đạo hàm đều mất. Với tính $C^1$ theo cả $x,t$, còn kiểm tra $u_t$. Với $C^2$, xét tiếp các góc của $g$ và đạo hàm cấp hai của $f$.

\textbf{Bước 5 — Tính thời điểm đến, giới hạn và ngưỡng.} Sóng ở $(x,t)$ chỉ nhận dữ liệu trong $[x-at,x+at]$. Giải bất đẳng thức giao hỗ trợ để tìm thời điểm đến và đi qua. Nếu $H$ hằng ngoài đoạn hỗ trợ, dùng các giá trị hằng này để tính giới hạn khi $t\to\infty$. Nếu hỏi áp suất cực đại, viết hàm từng khoảng thời gian và so sánh cực trị, bước nhảy cùng hai giá trị một phía.
\subsection*{Lời giải áp dụng cho bài}
\textbf{Bước 1. Chia dữ liệu và hai sóng.} $a=1$, $f(s)=s(s+1)$ trên $[-1,0]$, bằng $0$ ngoài đoạn, $g=0$. Vì vậy $A=B=f/2$, không cần nguyên hàm vận tốc.

\textbf{Bước 2. Dịch từng đoạn.} Sóng phải $A(x-t)$ khác không khi $t-1<x<t$, sóng trái $B(x+t)$ khác không khi $-t-1<x<-t$. Nghiệm chính xác là
\[u(x,t)=\tfrac12(x-t)(x-t+1)\chi_{[-1,0]}(x-t)+\tfrac12(x+t)(x+t+1)\chi_{[-1,0]}(x+t).\]
Khi hai đoạn chồng nhau, cộng hai đa thức; khi chỉ một đoạn hiện diện lấy một hạng; ngoài cả hai đoạn lấy $0$. Đây cũng là cách dựng các hình trong sách.

\textbf{Bước 3. Xét tính trơn.} $f$ bằng $0$ tại hai đầu nên $u$ liên tục. $f'$ trong đoạn bằng $2s+1$, có bước nhảy tại $-1,0$, nên các đường gãy cấp một là $x\pm t=-1,0$. Ở giao điểm $(-1/2,1/2)$, lát cắt $u(x,1/2)$ có đạo hàm trái theo $x$ bằng $1/2$, đạo hàm phải bằng $-1/2$; vì vậy điểm ấy vẫn không $C^1$. Tại giao của các mặt sóng phải so trực tiếp các đạo hàm một phía, không mặc nhiên giả sử chúng triệt tiêu.

\textbf{Bước 4. Kiểm tra.} Thế $t=0$ cho $u=f$, còn đạo hàm theo $t$ của hai sóng đối nhau nên vận tốc đầu bằng $0$ nơi $f'$ tồn tại. Mỗi đa thức của $x\pm t$ thỏa PDE; toàn nghiệm đúng theo nghĩa yếu trên các đường gãy.

\textbf{Chi tiết đạo hàm tại giao điểm hai mặt sóng.} Với $t=1/2$, bên trái $x=-1/2$, chỉ $x+t=x+1/2\in[-1,0]$ góp:
\[u(x,1/2)=\tfrac12(x+1/2)(x+3/2),\quad u_x=x+1\to1/2.\]
Bên phải, chỉ $x-t=x-1/2\in[-1,0]$ góp:
\[u(x,1/2)=\tfrac12(x-1/2)(x+1/2),\quad u_x=x\to-1/2.\]
Hai vết giá trị đều $0$, nhưng hai đạo hàm không bằng nhau, nên không thể gọi điểm giao là điểm trơn.\end{document}